% Fig. 2.12 / spec Fig A12 -- the mode as a corner solution: the density is
%   largest at the left end of its support, so the mode sits on the boundary.
% Source: mathstatch2.tex lines 2068-2103 (the first of the pair the
%   manuscript introduces with "two pictures to show the difference between a
%   corner solution and an interior one").
%
% Deviations from CH2_FIGURE_SPECS.md:
%
% 1. The manuscript curve 2.5 e^{-x} is a schematic: it does not integrate to
%    1.  Per spec 1.6 and spec A12's own note the web uses the true
%    Exponential(rate 1) density f(x) = e^{-x}, whose peak is f(0) = 1, and
%    marks that boundary point with a solid dot so that "the largest value is
%    reached at the edge" can be seen rather than inferred.
%
% 2. The word for "corner solution" is set below the picture in the
%    manuscript.  Spec 1.4 moves it to the caption; nothing Chinese is left
%    inside the drawing.
%
% 3. The manuscript puts the x label at (6.5,0.15) while the axis stops at
%    6.2.  Per spec 1.3 the label goes at the end of the axis instead.
%
% 4. \DeclareMathSizes{10}{10}{9}{8} raises the math script sizes from the
%    default 7pt/5pt so that the o of m_o and the X of f_X clear the 12px
%    floor of spec 1.3.  Same declaration as quartiles-equal-areas.tex; this
%    is not the site-wide \sssig change that ruling 21 defers.
%
% 5. The picture shares one explicit bounding box (the \path below) with
%    mode-interior-solution.tex, so that the two figures of the pair render at
%    exactly the same scale and their peaks, 2.9cm and 2.894cm, can be
%    compared directly.  Native size 220.402 x 111.135 pt, drawing
%    range 7.63cm wide, inside the 9.0cm cap; at the 350px phone width of P210
%    that puts 10pt text at 15.8px, 9pt scripts at 14.2px and 8pt
%    scriptscripts at 12.7px, and at 620px 28.0 / 25.2 / 22.4px.
\documentclass[tikz,border=2pt]{standalone}
\usepackage{amsmath}
\usepackage{amssymb}
\usepackage{xcolor}
\usetikzlibrary{arrows.meta}
\newcommand{\sssig}{\scriptscriptstyle}
\DeclareMathSizes{10}{10}{9}{8}

\definecolor{journalink}{HTML}{1F1C18}
\definecolor{journalmuted}{HTML}{8A7F72}
\definecolor{journalaccent}{HTML}{9C2D2D}

\begin{document}
\begin{tikzpicture}[
  x=1cm,
  y=2.9cm,
  every node/.style={font=\normalsize, journalink},
  axis/.style={journalink, line width=0.8pt, -{Stealth[length=4pt]}},
  tick/.style={journalink, line width=0.8pt},
  curve/.style={journalink, line width=0.8pt},
  guide/.style={journalmuted, line width=0.5pt, dash pattern=on 2pt off 2pt},
  solidpt/.style={circle, fill=journalink, inner sep=0pt, minimum size=3.6pt}
]
  % ---- bounding box shared with mode-interior-solution.tex --------------
  \path (-0.53cm,-0.80cm) rectangle (7.10cm,2.98cm);

  % ---- x axis ----------------------------------------------------------
  \draw[axis] (-0.5,0) -- (6.7,0);
  \node[anchor=west, inner sep=2.8pt] at (6.7,0) {$x$};

  % ---- the mode, here the left end point of the support ----------------
  \draw[guide] (0,1) -- (0,0);
  \draw[tick] ([yshift=0.035cm]0,0) -- ([yshift=-0.106cm]0,0);
  \node[anchor=north] at ([yshift=-0.25cm]0,0) {$m_{o}$};

  % ---- density curve: Exponential(1) -----------------------------------
  \draw[curve] plot[domain=0:6.4, samples=300, smooth] (\x,{exp(-\x)});
  \node[solidpt] at (0,1) {};

  % ---- curve label -----------------------------------------------------
  \node[anchor=south west, inner sep=0pt] at ([yshift=0.85cm]1.55,0)
    {$f_{\sssig X}(x)$};
\end{tikzpicture}
\end{document}
